\documentclass[10pt,a4paper]{amsart}
\usepackage{amsmath,amssymb,amsthm,graphicx}
% Retained minimal preamble for the Stacks Project chapter more-morphisms.tex (More on Morphisms).
% Verbatim excerpt of the chapter preamble: Xy-pic setup (native preamble lines 14-18) and the native
% theorem declarations with the mathematical macros (native preamble lines 154-217). Copied mechanically;
% nothing added, removed or redefined. The wrapper class and the amsmath/amssymb/amsthm/graphicx base
% packages are supplied by the assembler, so the class/externaldocument part of the native preamble is omitted.
% For commutative diagrams we use Xy-pic
\usepackage[all]{xy}

% We use 2cell for 2-commutative diagrams.
\xyoption{2cell}
% Theorem environments.
%
\theoremstyle{plain}
\newtheorem{theorem}[subsection]{Theorem}
\newtheorem{proposition}[subsection]{Proposition}
\newtheorem{lemma}[subsection]{Lemma}

\theoremstyle{definition}
\newtheorem{definition}[subsection]{Definition}
\newtheorem{example}[subsection]{Example}
\newtheorem{exercise}[subsection]{Exercise}
\newtheorem{situation}[subsection]{Situation}

\theoremstyle{remark}
\newtheorem{remark}[subsection]{Remark}
\newtheorem{remarks}[subsection]{Remarks}

\numberwithin{equation}{subsection}

% Macros
%
\def\lim{\mathop{\mathrm{lim}}\nolimits}
\def\colim{\mathop{\mathrm{colim}}\nolimits}
\def\Spec{\mathop{\mathrm{Spec}}}
\def\Hom{\mathop{\mathrm{Hom}}\nolimits}
\def\Ext{\mathop{\mathrm{Ext}}\nolimits}
\def\SheafHom{\mathop{\mathcal{H}\!\mathit{om}}\nolimits}
\def\SheafExt{\mathop{\mathcal{E}\!\mathit{xt}}\nolimits}
\def\Sch{\mathit{Sch}}
\def\Mor{\mathop{\mathrm{Mor}}\nolimits}
\def\Ob{\mathop{\mathrm{Ob}}\nolimits}
\def\Sh{\mathop{\mathit{Sh}}\nolimits}
\def\NL{\mathop{N\!L}\nolimits}
\def\CH{\mathop{\mathrm{CH}}\nolimits}
\def\proetale{{pro\text{-}\acute{e}tale}}
\def\etale{{\acute{e}tale}}
\def\QCoh{\mathit{QCoh}}
\def\Ker{\mathop{\mathrm{Ker}}}
\def\Im{\mathop{\mathrm{Im}}}
\def\Coker{\mathop{\mathrm{Coker}}}
\def\Coim{\mathop{\mathrm{Coim}}}

% Boxtimes
%
\DeclareMathSymbol{\boxtimes}{\mathbin}{AMSa}{"02}

%
% Macros for moduli stacks/spaces
%
\def\QCohstack{\mathcal{QC}\!\mathit{oh}}
\def\Cohstack{\mathcal{C}\!\mathit{oh}}
\def\Spacesstack{\mathcal{S}\!\mathit{paces}}
\def\Quotfunctor{\mathrm{Quot}}
\def\Hilbfunctor{\mathrm{Hilb}}
\def\Curvesstack{\mathcal{C}\!\mathit{urves}}
\def\Polarizedstack{\mathcal{P}\!\mathit{olarized}}
\def\Complexesstack{\mathcal{C}\!\mathit{omplexes}}
% \Pic is the operator that assigns to X its picard group, usage \Pic(X)
% \Picardstack_{X/B} denotes the Picard stack of X over B
% \Picardfunctor_{X/B} denotes the Picard functor of X over B
\def\Pic{\mathop{\mathrm{Pic}}\nolimits}
\def\Picardstack{\mathcal{P}\!\mathit{ic}}
\def\Picardfunctor{\mathrm{Pic}}
\def\Deformationcategory{\mathcal{D}\!\mathit{ef}}
\newenvironment{reference}{}{}
\newenvironment{slogan}{}{}
\newenvironment{history}{}{}
\usepackage{cleveref}
\crefname{originallemma}{original lemma}{original lemmas}
\Crefname{originallemma}{Original lemma}{Original lemmas}
\makeatletter
\newlabel{limits-lemma-higher-direct-images-zero-above-dimension-fibre}{{0E7D}{source}{}{sourcepointer-1}{}}
\newlabel{limits-lemma-higher-direct-images-zero-above-dimension-fibre@cref}{{[originallemma][0][]0E7D}{[1][0][]source}{}{}{}}
\makeatother
\crefname{originallemma}{original lemma}{original lemmas}
\Crefname{originallemma}{Original lemma}{Original lemmas}
\makeatletter
\newlabel{limits-lemma-proper-top-cohomology-finite-type}{{0EX4}{source}{}{sourcepointer-2}{}}
\newlabel{limits-lemma-proper-top-cohomology-finite-type@cref}{{[originallemma][0][]0EX4}{[1][0][]source}{}{}{}}
\makeatother
\crefname{originallemma}{original lemma}{original lemmas}
\Crefname{originallemma}{Original lemma}{Original lemmas}
\makeatletter
\newlabel{morphisms-lemma-openness-bounded-dimension-fibres}{{02FZ}{source}{}{sourcepointer-3}{}}
\newlabel{morphisms-lemma-openness-bounded-dimension-fibres@cref}{{[originallemma][0][]02FZ}{[1][0][]source}{}{}{}}
\makeatother
\title{Vanishing of H\textasciicircum{}1 of the structure sheaf on one curve fibre of a proper map spreads to a neighbourhood, universally}
\author{The Stacks Project Authors}
\date{stacks-project:more-morphisms.tex:local-snapshot-20261011; GNU FDL 1.2 or later}
\begin{document}
\maketitle
\noindent\textit{Editorial note.} Statement span is lines 20942-20949 of sources/stacks-more-morphisms/pinned-original-source.tex; proof span is lines 20951-20968, the complete author proof (16 body lines).\par\medskip
\noindent\textit{Editorial note.} Navigation stubs: 3 cited-but-not-included labels are supplied as non-load-bearing pointers with their pinned tags where the snapshot tag index has them.\par\medskip
\noindent\textit{Editorial note.} extra\_packages is empty on purpose: the retained preamble already loads xy with options and a second declaration clashes in the assembler wrapper.\par\medskip
\noindent\textit{Editorial note.} Text is retained verbatim, including the author's own typos and placeholders.\par\medskip
\noindent\textit{Editorial note.} This is the input lemma cited inside the collected card L16073 (globally generated vanishing); registered here as its own proof.\par\medskip
\noindent\textit{Allowed standard prerequisites (editorial).}\begin{itemize}
\item Explicit prerequisite boundary. Upper semicontinuity input for fibre dimension (Morphisms, lemma-openness-bounded-dimension-fibres) and the proper cohomology finiteness plus base-change theorems (Limits, lemmas proper-top-cohomology-finite-type, higher-direct-images-zero-above-dimension-fibre).
\end{itemize}\medskip
\ifdefined\hypertarget\hypertarget{sourcepointer-1}{}\fi
\noindent\textit{Source reference.} Original lemma 0E7D: sources/stacks-more-morphisms/pinned-original-source.tex, source lines 20968--20968. Cited by the retained statement or proof but not part of the collected spans; supplied only as a navigation pointer. The argument does not depend on its content.\par\medskip
\ifdefined\hypertarget\hypertarget{sourcepointer-2}{}\fi
\noindent\textit{Source reference.} Original lemma 0EX4: sources/stacks-more-morphisms/pinned-original-source.tex, source lines 20968--20968. Cited by the retained statement or proof but not part of the collected spans; supplied only as a navigation pointer. The argument does not depend on its content.\par\medskip
\ifdefined\hypertarget\hypertarget{sourcepointer-3}{}\fi
\noindent\textit{Source reference.} Original lemma 02FZ: sources/stacks-more-morphisms/pinned-original-source.tex, source lines 20968--20968. Cited by the retained statement or proof but not part of the collected spans; supplied only as a navigation pointer. The argument does not depend on its content.\par\medskip
\begin{lemma}
\label{lemma-h1-fibre-zero}
Let $f : X \to Y$ be a proper morphism of schemes. Let $y \in Y$
be a point with $\dim(X_y) \leq 1$ and $H^1(X_y, \mathcal{O}_{X_y}) = 0$.
Then there is an open neighbourhood $V \subset Y$ of $y$ such that
$R^1f_*\mathcal{O}_X|_V = 0$
and the same is true after base change by any $Y' \to V$.
\end{lemma}

\begin{proof}
To prove the lemma we may replace $Y$ by an open neighbourhood of $y$.
Thus we may assume $Y$ is affine and that all fibres of $f$ have
dimension $\leq 1$, see
Morphisms, Lemma \ref{morphisms-lemma-openness-bounded-dimension-fibres}.
In this case $R^1f_*\mathcal{O}_X$ is a quasi-coherent $\mathcal{O}_Y$-module
of finite type and its formation commutes with arbitrary base change, see
Limits, Lemmas \ref{limits-lemma-proper-top-cohomology-finite-type} and
\ref{limits-lemma-higher-direct-images-zero-above-dimension-fibre}.
Say $Y = \Spec(A)$, $y$ corresponds to the prime $\mathfrak p \subset A$, and
$R^1f_*\mathcal{O}_X$ corresponds to the finite $A$-module $M$.
Then $H^1(X_y, \mathcal{O}_{X_y}) = 0$ means that
$\mathfrak pM_\mathfrak p = M_\mathfrak p$ by the statement
on base change. By Nakayama's lemma
we conclude $M_\mathfrak p = 0$. Since $M$ is finite, we find
an $f \in A$, $f \not \in \mathfrak p$ such that $M_f = 0$.
Thus taking $V$ the principal open $D(f)$ we obtain the desired result.
\end{proof}

\end{document}
